\section{The Interior: Recursion, Integration, Selves}
\label{sec:interior}

The interior is the wild side of the substrate, but its structure is measurable, and the testbed
measures it. This section walks the findings on recursion and selves (P43, P57 to P67, P75). The
felt question, whether the structure of experience is the same as experiencing, is the open hard
problem named in the conceptual paper and not settled here.

\paragraph{Recursion without a new layer (P57, P58).}
A coherent, integrated cluster of vibes reads as a single higher vibe. P57 and P58 establish that
this higher vibe has no separately stored tone: it is the aggregate of its micro-tones, generalised
into a stable pattern and read as a unit, so the bottom stays ternary. The macro-rule is the same
signed-majority rule, recovered as a renormalization fixed point when the couplings are coarse-grained
along the system's genuine coherent regions. The model stays monist all the way up.

\paragraph{Nested selves and the tower (P59, P60).}
P59 treats a body made of cells: a small wound heals (homeostasis), while a whole cohesive part
flipped persists as a new identity (autonomy), with the crossover set by an integration threshold,
and the rest of the body left undisturbed. P60 builds the full tower, vibes into cells into tissues
into organs into a body, the same part-to-whole relation at every scale.

\paragraph{No infinite mirror (P61).}
P61 establishes that no part can hold a faithful copy of the whole, since a faithful copy needs at
least as many elements as the original. Self-representation is therefore lossy, a compressed summary,
and the endless tower of self-models converges within a finite total. The world knows itself only in
summary, which is why it can know itself at all.

\paragraph{Integration picks out selves (P63).}
P63 measures integration as the algebraic connectivity (the Fiedler value) of a region: a cohesive
cell has high integration, a random bag near zero, and a stable self sits at a local maximum. One
quantity governs whether a cluster is a being, whether it obeys the emergent rule, and whether a
disturbance is absorbed or kept.

\paragraph{Selves as attractors (P43, P75).}
P43 gives the structural account of freedom: a self is an attractor, a choice is that attractor
resolving an urge, determined yet self-determined and irreducible. P75 studies the attractors fully:
a self recovers from perturbations within a real basin (up to a forty percent flip), keeps its
identity over many beats (overlap $1.000$), and a mesh holds several selves at once, with the
capacity growing as the mesh grows.

\paragraph{Urge, dreaming, persistence, synchronicity (P64 to P67).}
P64 shows a slow, deep layer steering the fast surface as a real urge. P65 models waking and
dreaming as one mesh in two regimes, clamped to the world or free to roam its own stored patterns.
P66 shows pattern persistence, a self surviving a complete turnover of its material. P67 models
synchronicity as correlation between subsystems that share a deep common ancestry, the same
shared-substrate mechanism as the Bell correlations. These are structural models, with the felt
question held separate.
